\begin{titlepage}
\thispagestyle{empty}
\centering
\vspace*{1.6cm}

{\small\scshape\color{mutedgray} Technical Monograph \textbullet\ Relational Superposition Theory \par}
\vspace{0.9cm}

{\Huge\bfseries\color{inkblack} The Rank-Role Uncertainty Principle\par}
\vspace{0.45cm}
{\Large\color{slateink} A Formal Theory of Relational Superposition\par}

\vspace{1.2cm}
{\color{rulegray}\rule{0.35\textwidth}{0.6pt}}
\vspace{1.0cm}

{\large\bfseries Ratnadwip Sarkar\par}
\vspace{0.25cm}
{\small\color{slateink}
\url{https://github.com/RealRatnadwip} \quad\textbar\quad \url{https://ratnadwip.com}
\par}

\vspace{0.9cm}
{\small\color{mutedgray} Version 1.0 \quad\textbullet\quad \today\par}

\vspace{1.5cm}
\begin{minipage}{0.90\textwidth}
\small
\noindent\textbf{Abstract.}\;
We develop \emph{Relational Superposition Theory} (\RST), a mathematical framework for
systems in which a single focal individual maintains several concurrent romantic or
quasi-romantic connections, each fulfilling a distinct functional role (primary partner,
emotional support, financial support, entertainment, and contingency), without any
partner holding a well-defined ordinal rank. Prior to observation, no partner is first
or second; each partner occupies a weighted combination of roles and contributes a
probabilistic weight to the system state. We formalize rank and role as non-commuting
operators on a finite-dimensional inner-product space and establish the central
bound, the \emph{Rank-Role Inequality}: $\Delta R \cdot \Delta L \ge \kappa$, where
$\kappa > 0$ is the situationship constant governed by the commutator expectation.
We model the \emph{poke}, any direct status inquiry that forces projective collapse
onto a single partner, analyze the three post-poke phases of silence, weight
redistribution, and social-graph cascade (decoherence), and formalize the
\emph{Attractiveness-Optionality Conjecture} alongside its gender-symmetric mirror
configuration. Synthetic simulations and machine-checked Lean~4 proofs establish the
internal consistency and quantitative behavior of the framework.
\end{minipage}

\vspace{1.0cm}
\begin{minipage}{0.90\textwidth}
\small
\noindent\textbf{Keywords:} relational superposition, situationship, rank, role,
uncertainty, collapse, decoherence, optionality, social-graph cascade.
\end{minipage}

\vfill
\end{titlepage}

\clearpage
\pagenumbering{roman}
\tableofcontents
\clearpage
\listoffigures
\listoftables
\clearpage
\pagenumbering{arabic}

